%! Setting Up Determinants --- the 2 by 2 Determinant, Area, and Orientation — Unit 5, Lesson 5.0 — Warm-Up (Answer Key)
\documentclass[10pt]{article}
\usepackage{linalg-article}
\usepackage{linalg-key}   % pulls in -boxes; do NOT also load -boxes
\newcommand{\TallMath}[1]{$\displaystyle #1\rule[-1.4em]{0pt}{3.2em}$}

\begin{document}

\pageheader{Unit 5, Lesson 5.0}{Warm-Up --- Answer Key}
\namedateperiod

\vspace{0.12in}

\begin{notesbox}{Getting Ready: Ideas We'll Use Today}
\small These three ideas from Units~1 and~2 set up today's determinant. Answer quickly.

\vspace{4pt}
\textbf{1. Columns as vectors (Unit 1).} A matrix's columns are just vectors. Write the two
\emph{columns} of $A=\begin{bmatrix} 3 & 1 \\ 0 & 2 \end{bmatrix}$ as vectors:
\[
  \text{column 1} = \begin{bmatrix}{\color{keyred}\mathbf{3}}\\[2pt]{\color{keyred}\mathbf{0}}\end{bmatrix},
  \qquad
  \text{column 2} = \begin{bmatrix}{\color{keyred}\mathbf{1}}\\[2pt]{\color{keyred}\mathbf{2}}\end{bmatrix}.
\]

\vspace{4pt}
\textbf{2. Area of a box.} A rectangle has width $3$ and height $2$. Its area is
\[
  \text{area} = \text{base}\times\text{height} = {\color{keyred}\mathbf{6}}.
\]

\vspace{4pt}
\textbf{3. The inverse's denominator (Unit 2).} Recall
$A^{-1}=\dfrac{1}{ad-bc}\begin{bmatrix} d & -b \\ -c & a \end{bmatrix}$. For
$A=\begin{bmatrix} 2 & 1 \\ 1 & 3 \end{bmatrix}$, compute the number $ad-bc$ in the denominator, and
say when the inverse would \emph{fail} to exist.
\ansline{$ad-bc=(2)(3)-(1)(1)=6-1=5$. The inverse fails when $ad-bc=0$ (you cannot divide by $0$).}
\end{notesbox}

\begin{teachernote}
All three items set up the determinant. Item~1 reminds students the columns are vectors --- today
those vectors span a parallelogram. Item~2 primes ``area $=$ base $\times$ height.'' Item~3 is the
hinge: the number $ad-bc$ they already divide by in the Unit~2 inverse \emph{is} today's determinant,
and ``the inverse dies when $ad-bc=0$'' is the fact the lesson turns on. Watch for anyone writing
$ad+bc$ in item~3.
\end{teachernote}

\end{document}
